\chapter{Valutazione finale}
\label{cap:valutazione_finale}

\section{Obiettivi raggiunti}
Il progetto ha raggiunto i seguenti obiettivi primari:

\begin{itemize}
    \item \textbf{Implementazione completa del protocollo OID4VP\glsadd{oid4vpg}\glsfirstoccur{}:} il sistema supporta l'intero ciclo di vita delle \textit{Verifiable Presentations}, dalla richiesta iniziale fino al ritiro dei dati verificati, coprendo le quattro fasi previste dalla specifica.
    \item \textbf{Supporto a sei schemi di identificazione del client:} il sistema gestisce correttamente \texttt{redirect\_uri}, \texttt{x509\_san\_dns}, \texttt{x509\_hash}, \texttt{openid\_federation}, \texttt{verifier\_attestation} e \texttt{pre\_registered}, coprendo la quasi totalità degli scenari previsti dal draft OID4VP\glsadd{oid4vpg}.
    \item \textbf{Verifica crittografica SD-JWT\glsadd{sdjwtg}\glsfirstoccur{} end-to-end:} il sistema è in grado di analizzare, validare gli hash delle \textit{disclosure}, verificare la firma dell'Issuer e convalidare il \textit{Key Binding JWT\glsadd{jwtg}\glsfirstoccur{}}, garantendo sia l'autenticità sia la titolarità delle credenziali.
    \item \textbf{Risoluzione ricorsiva delle Trust Chain\glsadd{trustchaing}\glsfirstoccur{}:} l'algoritmo DFS\glsadd{dfs}\glsfirstoccur{} implementato nel \texttt{TrustAnchorService} è in grado di risolvere catene di fiducia arbitrariamente profonde con protezione contro i cicli.
    \item \textbf{Sicurezza Forward Secrecy\glsadd{forwardsecrecyg}\glsfirstoccur{}:} il meccanismo di chiavi effimere con distruzione atomica tramite \texttt{GETDEL} garantisce che la compromissione di una chiave non comprometta sessioni passate.
    \item \textbf{Multi-tenancy\glsadd{multitenancyg}\glsfirstoccur{} nativo:} l'isolamento dei dati per dominio è garantito sia a livello applicativo sia a livello di database, mediante indici composti su MongoDB.
    \item \textbf{Superamento dei Conformance Test OID4VP\glsadd{oid4vpg}:} il sistema ha superato con successo i test di conformità ufficiali della OpenID Foundation per il profilo Verifier, come descritto nel Capitolo~\ref{cap:verifica_validazione}, certificando l'interoperabilità e la correttezza del comportamento di protocollo.
\end{itemize}

\section{Debito tecnico e limiti}
È necessario riconoscere con rigore i limiti del lavoro svolto, che costituiscono il debito tecnico del progetto:

\begin{enumerate}
    \item \textbf{Assenza di test automatizzati.} Il progetto non include test unitari, di integrazione né end-to-end. Non è installato alcun framework di testing (Jest, Vitest, Mocha), nonostante la presenza di \texttt{ioredis-mock} nelle \texttt{devDependencies} e la riga commentata \texttt{\# RUN yarn test} nel \texttt{Dockerfile} suggeriscano che il testing fosse previsto. Questa rappresenta la criticità più rilevante: la verifica dinamica dipende interamente dal testing manuale, intrinsecamente non ripetibile e non scalabile, e non è possibile integrare la verifica in una pipeline CI/CD\glsadd{cicdg}\glsfirstoccur{} con gate automatici.
    \item \textbf{Curva EC limitata a P-256.} L'\texttt{EphemeralKeyService} genera sempre chiavi effimere con curva P-256, indipendentemente dall'algoritmo JWE\glsadd{jweg}\glsfirstoccur{} negoziato. Algoritmi come \texttt{ECDH-ES+A256KW} (che beneficerebbe di P-384 o P-521) e lo scambio basato su X25519 non sono pertanto supportati.
    \item \textbf{Disclosure di tipo array non supportate.} Il \texttt{SdJwtCredentialVerifierService} gestisce esclusivamente le \textit{disclosure} a tre elementi (salt, nome del claim, valore). Le \textit{disclosure} a due elementi (salt, valore), utilizzate per elementi all'interno di array, producono esplicitamente un errore. Questo limita la capacità del sistema di gestire credenziali SD-JWT\glsadd{sdjwtg} con claim a valori multipli.
    \item \textbf{Assenza di caching per le Entity Configuration.} Ogni invocazione del \texttt{TrustAnchorService} per la validazione di una Trust Chain\glsadd{trustchaing} esegue chiamate HTTP\glsadd{http}\glsfirstoccur{} verso gli endpoint \texttt{/.well-known/openid-federation} senza alcun meccanismo di cache. In scenari con molteplici richieste per il medesimo Issuer, questo genera traffico di rete ridondante e latenza evitabile.
    \item \textbf{Verifica della firma dell'Attestation JWT\glsadd{jwtg} mancante.} Nel metodo \texttt{validateVerifierAttestation} del \texttt{ClientIdPrefixValidator}, il JWT\glsadd{jwtg} di attestazione viene decodificato per estrarne i claim, ma la firma non viene verificata crittograficamente contro la chiave pubblica dell'autorità emittente.
    \item \textbf{Copertura parziale dei formati di credenziali.} I Conformance Test eseguiti hanno coperto il formato \texttt{dc+sd-jwt}, ma il sistema dichiara di poter supportare anche il formato \texttt{mso\_mdoc} (ISO 18013-5) tramite la struttura \texttt{MdocFormatSupported} nel modello \texttt{Client}. Nessun servizio di verifica per questo formato è stato implementato, e la discrepanza tra le capacità dichiarate e quelle testate costituisce un rischio di conformità.
    \item \textbf{Assenza di test di performance e carico.} Non sono stati condotti test di performance, load testing o stress testing. L'assenza di dati sulle performance sotto carico (latenza media, throughput, comportamento sotto contention su Redis) rappresenta un'incognita per il dimensionamento dell'infrastruttura in scenari di produzione.
    \item \textbf{Validazione limitata ai flussi di protocollo.} I Conformance Test validano la conformità del protocollo dall'esterno, ma non entrano nel merito della correttezza interna dell'implementazione: non verificano, ad esempio, che la chiave privata effimera venga effettivamente distrutta dopo l'uso (\textit{Forward Secrecy}\glsadd{forwardsecrecyg}), che il consumo della sessione sia realmente atomico sotto concorrenza, o che il convertitore di nomenclatura non corrompa strutture dati con chiavi non previste.
\end{enumerate}

\section{Sviluppi futuri e miglioramenti proposti}
Sulla base dell'analisi dei limiti, i principali assi di miglioramento per le iterazioni successive sono:

\begin{itemize}
    \item \textbf{Introduzione di una suite di test completa:} installazione di Jest o Vitest, scrittura di test unitari per i servizi di validazione crittografica (utilizzando \texttt{ioredis-mock} per Redis) e test di integrazione per i flussi end-to-end. La priorità dovrebbe essere assegnata ai servizi \texttt{trust-anchor.service.ts}, \texttt{credential-verifier.service.ts} e \texttt{prefix-validator.service.ts}, dove un errore è più critico.
    \item \textbf{Integrazione nella pipeline CI/CD\glsadd{cicdg}:} decommentare la riga \texttt{RUN yarn test} nel \texttt{Dockerfile} e configurare una pipeline di CI/CD (ad esempio GitLab CI o GitHub Actions) che esegua automaticamente \texttt{yarn typecheck}, \texttt{yarn lint}, \texttt{yarn test} e \texttt{yarn audit} ad ogni commit, bloccando il merge in caso di fallimento.
    \item \textbf{Test di integrazione end-to-end e di concorrenza:} scrittura di test che simulino l'intero flusso delle quattro fasi OID4VP\glsadd{oid4vpg} in un ambiente controllato, verificando le transizioni di stato, la generazione e distruzione delle chiavi effimere, e la correttezza dei payload. Inoltre, test specifici che invochino \texttt{consumeSessionForResponse()} e \texttt{retrieveAndDestroyPrivateKey()} in modo concorrente per verificare che l'atomicità sia garantita sotto contention.
    \item \textbf{Supporto a curve EC multiple e X25519:} parametrizzazione della curva EC nell'\texttt{EphemeralKeyService} in base all'algoritmo JWE\glsadd{jweg} negoziato, con supporto a P-384, P-521 e, in prospettiva, X25519 per lo scambio chiavi Diffie-Hellman.
    \item \textbf{Implementazione delle array \textit{disclosure}:} completamento del supporto a \textit{disclosure} a due elementi nel \texttt{SdJwtCredentialVerifierService}, abilitando la verifica di credenziali SD-JWT\glsadd{sdjwtg} con claim a valori multipli.
    \item \textbf{Implementazione di un layer di caching per le Entity Configuration:} salvataggio in Redis (con TTL\glsadd{ttlg}\glsfirstoccur{} basato sul claim \texttt{exp}) delle \textit{Entity Configuration} e dei \textit{Subordinate Statement} recuperati durante la risoluzione delle Trust Chain\glsadd{trustchaing}, al fine di ridurre la latenza e il traffico di rete.
    \item \textbf{Verifica crittografica completa dell'Attestation JWT\glsadd{jwtg}:} implementazione della verifica della firma nel metodo \texttt{validateVerifierAttestation}, importando la chiave pubblica dell'autorità emittente e validando il JWT\glsadd{jwtg} con \texttt{jose.jwtVerify()}.
    \item \textbf{Estensione dei Conformance Test:} una volta implementato il supporto al formato mDoc, eseguire i Conformance Test OIDF anche per questo formato. Valutare inoltre l'esecuzione dei test per il profilo HAIP\glsadd{haipg}\glsfirstoccur{} (\textit{High Assurance Interoperability Profile}), che include requisiti aggiuntivi di sicurezza rilevanti per il contesto EUDI.
    \item \textbf{Rafforzamento della configurazione TypeScript:} l'attivazione di \texttt{strictNullChecks} (o dell'intero \texttt{strict: true}) forzerebbe il compilatore a segnalare staticamente tutti i punti in cui un valore potrebbe essere null, eliminando un'intera categoria di errori a runtime.
\end{itemize}